% Source: https://artofproblemsolving.com/community/c3h1849670_equation_of_circle
% Stable identity: aops-topic-1849670-q01
The radius is $\dfrac{10}{2}=5$. A point $(x,y)$ is on the circle exactly when its distance from $(-2,3)$ is $5$:
\[
\sqrt{(x+2)^2+(y-3)^2}=5 .
\]
Squaring gives
\[
\boxed{(x+2)^2+(y-3)^2=25},
\]
or, expanded, $x^2+y^2+4x-6y-12=0$.

\begin{asy}
usepackage("amsmath");
size(5.5cm);
defaultpen(fontsize(10pt));
pair O=(-2,3);
draw((-8,0)--(4.2,0),Arrow(TeXHead));draw((0,-2.6)--(0,8.8),Arrow(TeXHead));
label("$x$",(4.2,0),S);label("$y$",(0,8.8),W);
draw(circle(O,5),blue+linewidth(1));
draw((-7,3)--(3,3),dashed);
dot(O);dot((-7,3));dot((3,3));
label("$(-2,3)$",O,N);label("$10$",(-4.5,3),S);
label("$(-7,3)$",(-7,3),W);label("$(3,3)$",(3,3),E);
label("\mbox{}",(min(currentpicture,user=true)+max(currentpicture,user=true))/2,black); // keeps the text after the figure black when LaTeX typesets the labels inline
\end{asy}

For example, the endpoints $(-7,3)$ and $(3,3)$ of the horizontal diameter satisfy the equation: $(-5)^2+0^2=25$ and $5^2+0^2=25$.
